%==============================================================================
% PART 2: Methods Section (Fixed - No Duplicates)
%==============================================================================
\section{Methods}

\subsection{Pipeline Architecture}

GlycoSMILES2BAP operates through four sequential modules: Input Parsing, CCD Mapping, BAP Generation, and Output Formatting. The pipeline accepts IUPAC-condensed, WURCS, and GlycoCT input formats.

\subsection{CCD Mapper Module}

The CCD (Chemical Component Dictionary) provides standardized residue identifiers. Our mapper supports 28+ monosaccharide configurations. Key design decisions include: (1) case-insensitive matching, (2) anomeric position tracking (C2 for sialic acids, C1 for aldoses), and (3) ring oxygen positions (O4 for pentoses, O5 for hexoses, O6 for sialic acids).

\begin{table}[h]
\centering
\caption{Key CCD Mappings. The complete mapping table for 28+ monosaccharides is provided in Supplementary Table S1.}
\begin{tabular}{lllll}
\toprule
Monosaccharide & Anomer & Config & CCD Code & Anomeric C \\
\midrule
GlcNAc & $\beta$ & D & NAG & C1 \\
GlcNAc & $\alpha$ & D & A2G & C1 \\
Man & $\alpha$ & D & MAN & C1 \\
Man & $\beta$ & D & BMA & C1 \\
Gal & $\beta$ & D & GAL & C1 \\
Gal & $\alpha$ & D & GLA & C1 \\
Fuc & $\alpha$ & L & FUC & C1 \\
Neu5Ac & $\alpha$ & D & SIA & C2 \\
Neu5Gc & $\alpha$ & D & NGC & C2 \\
Xyl & $\beta$ & D & XYS & C1 \\
\bottomrule
\end{tabular}
\label{tab:ccd_mapping}
\end{table}

\subsection{BAP Generator Module}

The bondedAtomPairs (BAP) generator converts glycan topology into AF3-compatible bond specifications. Each glycosidic linkage is represented as an explicit atom pair specifying donor residue, donor atom, acceptor residue, and acceptor atom.

\textbf{Algorithm 1: BAP Generation}

\noindent\textbf{Input:} Glycan topology $T$, CCD mapper $M$\\
\textbf{Output:} List of bondedAtomPairs $B$

\begin{enumerate}
\item $B \gets []$
\item \textbf{for} each linkage $l$ in $T$.edges \textbf{do}
\item \quad $donor\_res \gets l$.donor\_id
\item \quad $acceptor\_res \gets l$.acceptor\_id
\item \quad $donor\_atom \gets M$.anomeric($T$.nodes[$donor\_res$].ccd)
\item \quad $acceptor\_atom \gets$ ``O'' + $l$.acceptor\_position
\item \quad $B$.append([$donor\_res$, $donor\_atom$, $acceptor\_res$, $acceptor\_atom$])
\item \textbf{end for}
\item \textbf{return} $B$
\end{enumerate}

\noindent The algorithm runs in $O(n)$ time where $n$ is the number of monosaccharide residues, making it efficient for large glycans.

\subsection{PDB Structure Validation Method}
\label{sec:pdb_validation}

To independently validate the CCD mapper's stereochemistry handling, we employed a PDB structure comparison approach. This validation was necessary because AlphaFold Server does not support direct glycan submission, and AF3 model parameters were not accessible at the time of this study.

We selected four PDB structures with known correct stereochemistry, each testing a distinct aspect of CCD mapping (Table~\ref{tab:pdb_cases}).

\begin{table}[h]
\centering
\caption{PDB Validation Test Cases}
\begin{tabular}{lllp{5cm}}
\toprule
ID & PDB & Structure & Key Validation Point \\
\midrule
V001 & 5NSC & Fucosylated & FUC = alpha-L-fucose (L-configuration) \\
V002 & 1L6R & Lactose & GAL C4 axial hydroxyl (epimer distinction) \\
V003 & 2VXR & Sialyllactose & SIA anomeric = C